% =====================================================================
%  CHAPTER 2 --- THEORETICAL BACKGROUND
% =====================================================================
\chapter{THEORETICAL BACKGROUND}
\label{ch:theory}

This chapter builds a single chain of reasoning, from the Sun to the
simulated image, and every equation below is one specific link in that
chain-not an independent piece of theory. Told as one sentence: the Sun
produces radio emission; the international community measures this emission
as a \emph{flux density}, in solar flux units; a receiving antenna converts
that flux density into received radio-frequency \emph{power}; a receiver
converts that power into a \emph{voltage}, together with its own unavoidable
thermal noise; several antennas at known positions see the same signal with a
small, geometry-dependent phase difference; a \emph{correlator} combines a
pair of antenna voltages into a single complex number called the
\emph{visibility}; the array geometry places every visibility at its correct
point in a Fourier $(u,v)$ plane; and an inverse Fourier transform of those
points produces the \emph{dirty image} analysed in Chapter~\ref{ch:results}.
The sections below develop this chain in order, introducing each equation
only once the preceding step has produced the quantity it needs. Not every
equation plays the same role, and this chapter says so explicitly at each
step: some equations (for example Eq.~\eqref{eq:friis}, \eqref{eq:vrms},
\eqref{eq:johnson}, and \eqref{eq:uvw}) are computational steps that the
simulation of Chapter~\ref{ch:pipeline} evaluates directly; others (for
example Eq.~\eqref{eq:fringeperiod}, \eqref{eq:vcz}-\eqref{eq:jinc}, and
\eqref{eq:dirty}) are analytic results that the simulation is not built from,
but is instead checked against, as validation benchmarks in
Chapter~\ref{ch:results}.

In outline: Section~\ref{sec:solar-activity} places the solar radio burst
inside the wider picture of solar activity and motivates the observing
frequency used in this thesis. Section~\ref{sec:radioastro-quantities} turns
that physical picture into the standard radio-astronomy quantities-brightness
temperature and flux density-that the simulation can actually work with.
Section~\ref{sec:flux-to-voltage} carries the flux density through the
receiving antenna to a voltage at the receiver input, and
Section~\ref{sec:noise} adds the unavoidable thermal noise and states how
averaging improves sensitivity. Section~\ref{sec:voltage-to-visibility}
explains how two such voltages are combined by a correlator into a single
complex visibility. Section~\ref{sec:uvw-coverage} shows how the physical
antenna positions and the Earth's rotation determine which visibilities can
actually be measured. Section~\ref{sec:visibility-to-image} closes the chain
by turning a set of measured visibilities back into an image of the sky.
Section~\ref{sec:chapter-summary} collects the chain into one summary and
sorts its equations into the two roles above. The treatment of interferometry
and synthesis imaging follows closely the formalism of
\citet{taylor1999synthesis}, \emph{Synthesis Imaging in Radio Astronomy II},
whose Chapters~1-3 are the principal reference for
Sections~\ref{sec:voltage-to-visibility}-\ref{sec:visibility-to-image}.
Equations are numbered within the chapter for cross-reference in the later
design and results chapters.

% =====================================================================
\section{Solar Radio Activity and Observation Target}
\label{sec:solar-activity}

\subsection{Solar Flare, CME, Geomagnetic Storm, and Solar Radio Burst}
\label{subsec:flare-cme-storm}
The Sun does not only shine in visible light. It radiates across almost the
entire electromagnetic spectrum, from X-rays to radio waves, and each part of
that spectrum reveals a different physical process. Chapter~\ref{ch:intro}
motivated this thesis from the practical side: solar activity can disturb
satellites, GNSS, and power grids, and this is why monitoring the Sun matters.
This section gives the physical picture behind that motivation, introducing
the events in the order in which they actually happen.

A \emph{solar flare} is a sudden release of magnetic energy stored in the
corona, usually above a sunspot group, which heats the local plasma and
accelerates electrons and ions to high energy on timescales of
minutes~\citep{dulk1985,mclean1985}. A \emph{coronal mass ejection} (CME) is a
related but distinct event: a large-scale eruption of coronal plasma and
magnetic field that is expelled outward, sometimes reaching the Earth in one
to three days. A flare or a CME-driven shock can also accelerate protons and
electrons to much higher energies than the ambient plasma; when these escape
into interplanetary space they are observed as a \emph{solar energetic
particle} (SEP) event. When a CME, a SEP event, or a fast stream of solar
wind reaches the Earth's magnetosphere, it can compress and disturb it,
producing a \emph{geomagnetic storm}-the event responsible for the
satellite-drag incident described in
Chapter~\ref{ch:intro}~\citep{berger2023starlink}. A geomagnetic storm, in
turn, can change the electron density of the Earth's upper atmosphere,
producing an \emph{ionospheric disturbance} that degrades GNSS signals, as
discussed in Chapter~\ref{ch:intro}~\citep{floressoriano2024gnss}.

A \emph{solar radio burst} (SRB) is the radio-frequency signature of the same
energetic processes: as electrons are accelerated and plasma is disturbed by
flares, shocks, and CMEs, they emit radio waves that are, unlike the CME
plasma itself, detectable from the ground within seconds. This is why solar
radio bursts are useful as an early, remote indicator of the eruptive
activity that eventually leads to a geomagnetic storm or an ionospheric
disturbance. Table~\ref{tab:solar-terms} collects the terms introduced above,
for reference throughout the rest of this chapter.

\begin{table}[H]
  \centering
  \caption[Key terms used to describe solar activity and solar radio emission]{Key terms used to describe solar activity and solar radio emission.}
  \label{tab:solar-terms}
  \renewcommand{\arraystretch}{1.3}
  \begin{tabularx}{\textwidth}{l >{\raggedright\arraybackslash}X}
    \toprule
    \textbf{Term} & \textbf{Meaning} \\
    \midrule
    Solar flare & Sudden release of magnetic energy in the corona; heats plasma and accelerates particles on a timescale of minutes. \\
    Coronal mass ejection (CME) & Large-scale eruption of coronal plasma and magnetic field, expelled outward from the Sun. \\
    Solar energetic particle (SEP) event & Burst of high-energy protons and electrons, accelerated by a flare or a CME-driven shock, that can reach the Earth. \\
    Geomagnetic storm & Disturbance of the Earth's magnetosphere caused by a CME, a SEP event, or a fast solar-wind stream; can increase atmospheric drag on satellites (Chapter~\ref{ch:intro}). \\
    Ionospheric disturbance & Change in the electron density of the Earth's ionosphere, often following a geomagnetic storm, that can degrade GNSS signals (Chapter~\ref{ch:intro}). \\
    Solar radio burst (SRB) & Radio-frequency signature of the energetic processes above; classified into Types~I-V (Section~\ref{subsec:burst-types}). \\
    Quiet Sun & Steady, thermal radio background present even when no burst is occurring (Section~\ref{subsec:quiet-sun-brightness}). \\
    \bottomrule
  \end{tabularx}
\end{table}

\subsection{Solar Radio Burst Types and Observing Frequency Ranges}
\label{subsec:burst-types}
At radio wavelengths the Sun is not the sharp-edged disk seen in visible light
but a stratified, optically thick plasma whose apparent size and brightness
depend strongly on the observing frequency~\citep{kundu1965,dulk1985}. Three
atmospheric layers matter to a radio observer: the photosphere (the
\SI{5800}{\kelvin} optical surface, dominant only at millimetre wavelengths),
the chromosphere, and the hot, tenuous corona, whose electron temperature
reaches $\sim\!\SI{e6}{\kelvin}$. At metre and decimetre wavelengths the
radiation received originates predominantly in the corona and upper
chromosphere.

Superimposed on the steady, thermal background (Section~\ref{subsec:quiet-sun-brightness})
are the transient solar radio bursts introduced above, which can raise the
flux by several orders of magnitude. Following the classical morphological
scheme, they are grouped into five types (I-V) according to their appearance
in the dynamic spectrum, and each type is associated with a characteristic
frequency range~\citep{wild1963,mclean1985}:
\begin{itemize}[leftmargin=1.5em]
  \item \textbf{Type~I} noise storms: long-lived, narrowband emission,
        typically \SIrange{50}{300}{\mega\hertz}.
  \item \textbf{Type~II} bursts: slowly drifting (of order \SI{1}{\mega\hertz\per\second}),
        associated with coronal and interplanetary shock waves, typically
        starting below \SI{150}{\mega\hertz} and drifting down to hundreds of
        kilohertz as the shock propagates outward.
  \item \textbf{Type~III} bursts: fast-drifting, produced by electron beams
        streaming along open field lines; they can start above
        \SI{1}{\giga\hertz} and sweep down through the decimetre and metre
        bands within seconds.
  \item \textbf{Type~IV} continua: broadband, following flares, spanning the
        decimetre to metre range and sometimes lasting hours.
  \item \textbf{Type~V} continua: short (tens of seconds), sometimes following
        Type~III bursts.
\end{itemize}
These frequency ranges span several named observing bands: \emph{decametric}
(below about \SI{30}{\mega\hertz}), \emph{metric} (roughly
\SIrange{30}{300}{\mega\hertz}), \emph{decimetric} or \emph{UHF} (roughly
\SI{300}{\mega\hertz} to \SI{3}{\giga\hertz}), and \emph{microwave} (above
\SI{3}{\giga\hertz}). The \SI{610}{\mega\hertz} band used in this thesis lies
inside the decimetric/UHF range. At decimetre wavelengths, the emission most
often seen is the fast-drifting tail of Type~III bursts as they sweep through
the band, the decimetric/microwave continuum associated with Type~IV events,
and narrowband spike bursts produced by non-thermal (plasma or
gyrosynchrotron) emission mechanisms. This emission is spatially compact,
which is precisely what makes a burst a meaningful target for a small
interferometer.

\begin{figure}[H]
  \centering
  \begin{overpic}[width=\textwidth]{fig_solar_dynamic_spectrum}
    \put(68,25){\color{white}\rule{22\unitlength}{4\unitlength}}
    \put(79,26.5){\makebox(0,0){\textcolor{blue!40!black}{\bfseries\large I~Storm}}}
  \end{overpic}
  \caption[Solar dynamic spectrum showing Type I--V burst classification]{Solar dynamic spectrum showing Type I--V burst classification.\\[4pt]
  The schematic frequency--time plot illustrates the quiet-Sun background
  continuum and the five burst categories. Type~I storms appear as
  narrowband chains; the fast-drifting Type~III lanes are produced by electron
  beams escaping along open field lines and are the dominant transient signature
  at decimetre wavelengths; slowly drifting Type~II arcs trace coronal shock
  fronts driven by coronal mass ejections.
  Adapted from \citet{mclean1985} and \citet{dulk1985}; the long-duration,
  narrowband Type~I noise-storm label on the source figure has been corrected
  here (originally mislabelled).}
  \label{fig:dynamic-spectrum}
\end{figure}

\subsection{Why 610 MHz, and Why the Quiet Sun Is Useful for Validation}
\label{subsec:why-610mhz}
Chapter~\ref{ch:intro} explained the practical reasons for working at
\SI{610}{\mega\hertz}: archived flux data are freely available at this
frequency, and low-cost UHF Yagi-Uda antennas are easy to obtain. The physical
picture above adds the remaining justification. First, \SI{610}{\mega\hertz}
still falls inside the decimetric/UHF band where flare-related non-thermal
emission-the fast-drifting tail of Type~III bursts, decimetric continua, and
spike bursts-is observed, so the band is not physically irrelevant to solar
radio-burst work, even though the strongest and most frequently studied
bursts (Type~II and Type~III at their onset) are usually observed at lower
frequencies, tens of MHz, closer to the metric/decametric band used by
instruments such as DLITE~\citep{carson2022dlite}. Second, and just as
important for this thesis, the \emph{quiet Sun}-the steady, thermal background
described in Section~\ref{subsec:quiet-sun-brightness}-is present at
\SI{610}{\mega\hertz} every day, regardless of whether a burst occurs. Its
brightness is stable and well characterised, so it provides a predictable,
always-available source against which the geometry, timing, and signal chain
of the simulated interferometer can be validated. A burst is unpredictable
and, for a first-generation instrument, this thesis therefore uses the quiet
Sun, observed during the transit of \mbox{20 March 2025} at Learmonth, as the
validation target: any departure of the simulated response from the expected
analytic result can then be attributed to the instrument model rather than to
unpredictable source variability.

% =====================================================================
\section{From Solar Radio Emission to Flux Density}
\label{sec:radioastro-quantities}

\subsection{Specific Intensity and Brightness Temperature}
\label{subsec:quiet-sun-brightness}
A radio telescope does not measure temperature directly, the way a
thermometer does. It measures the intensity of the radio radiation arriving
from the sky. In the absence of activity, the \emph{quiet Sun} radiates by
thermal free-free (bremsstrahlung) emission, and in radio astronomy it is
conventional to express this intensity as an equivalent \emph{brightness
temperature} $T_b$: the temperature that a blackbody would need to have in
order to produce the same observed specific intensity $I_\nu$ at frequency
$\nu$. In the radio (Rayleigh-Jeans) limit, this definition is written as
\begin{equation}
  I_\nu = \frac{2 k_B \nu^2}{c^2}\,T_b ,
  \label{eq:rj}
\end{equation}
where:
\begin{itemize}[leftmargin=1.5em,itemsep=1pt,topsep=2pt]
  \item $I_\nu$ is the specific intensity, in $\si{\watt\per\square\metre\per\hertz\per\steradian}$;
  \item $T_b$ is the brightness temperature, in kelvin;
  \item $\nu$ is the observing frequency, in hertz;
  \item $k_B$ is Boltzmann's constant, which appears because it is the
        physical constant that converts a temperature in kelvin into an
        energy scale;
  \item $c$ is the speed of light \citep{dulk1985,kraus1986}.
\end{itemize}
At \SI{610}{\mega\hertz} the quiet-Sun brightness temperature is of order
$10^5$-$10^6~\si{\kelvin}$. Equation~\eqref{eq:rj} is a \emph{definition} of
what brightness temperature means-it is not a computational step that the
simulation evaluates. The simulation never needs a numerical value of $T_b$:
it works directly with the flux density introduced next, which is the
quantity an antenna actually collects. Brightness temperature is given here
only because it is the standard way the radio-astronomy literature describes
how bright the quiet Sun is, and because it makes the stability argument of
Section~\ref{subsec:why-610mhz} precise: the quiet Sun's brightness
temperature is steady from day to day, so departures from the predicted
fringe pattern can be attributed to the instrument rather than to the source.

\subsection{Flux Density and the Solar Flux Unit}
\label{subsec:flux-sfu}
The specific intensity $I_\nu$ describes the brightness of a resolved source
per unit solid angle, which is not directly what a receiver measures. What a
receiver actually collects is the \emph{flux density} $S_\nu$-the intensity
integrated over the angular extent of the source-expressed in
$\si{\watt\per\square\metre\per\hertz}$. For the Sun, the integrated flux
density is a few tens of solar flux units \citep{dulk1985,kundu1965}, where
one solar flux unit is
\begin{equation}
  1~\sfu = \SI{e-22}{\watt\per\square\metre\per\hertz}.
  \label{eq:sfu}
\end{equation}
The solar flux unit is the standard scale on which archived solar radio
observatories, including Learmonth, report their measurements. In this
thesis, the archived Learmonth flux record $S_\nu(t)$, in solar flux units at
\SI{610}{\mega\hertz}, is the external physical input of the simulation: it
is why the solar flux unit matters here, and it is the quantity that
Sections~\ref{sec:flux-to-voltage}-\ref{sec:visibility-to-image} carry through
every remaining stage of the theoretical chain.

\subsection{From Flux Density to a Simulation Input \texorpdfstring{$S_\nu(t)$}{Sν(t)}}
\label{subsec:sv-of-t}
The Learmonth archive does not provide a continuous function of time; it
provides a sequence of flux-density samples recorded at fixed intervals over
the course of a day, in SFU, at \SI{610}{\mega\hertz}. To use this record as
the input of a time-domain simulation, the discrete samples are converted to
SI units using Eq.~\eqref{eq:sfu} and then interpolated to produce a
continuous function $S_\nu(t)$ that can be evaluated at every simulation time
step, however finely the array geometry or receiver chain is sampled. The
details of the interpolation, and the archived file format from which
$S_\nu(t)$ is built, are given in Chapter~\ref{ch:pipeline}. What matters here
is the role $S_\nu(t)$ plays in the theoretical chain: it is the single
external input that carries the real, measured behaviour of the Sun into
every later stage of the simulation, from the received voltage
(Section~\ref{sec:flux-to-voltage}) to the final image
(Section~\ref{sec:visibility-to-image}).

% =====================================================================
\section{From Solar Flux to Received Voltage}
\label{sec:flux-to-voltage}
A solar flux density is not yet an electrical signal. To simulate a
receiver, $S_\nu(t)$ must be converted, step by step, into a received
radio-frequency power and then into a voltage that a digitiser could
actually sample. This section carries out that conversion.

\subsection{Antenna Gain, Primary Beam, and Effective Aperture}
\label{subsec:antenna-gain-beam}
A receiving antenna converts the incident electromagnetic field into a
voltage at its terminals. Its directional sensitivity is described by the
\emph{normalised power pattern} $A_n(\theta,\phi)$, which for a receiving
element is identical to its transmitting pattern by reciprocity. The
\emph{directivity} $D$ and \emph{gain} $G$ follow from
\begin{equation}
  D = \frac{4\pi}{\displaystyle\iint_{4\pi} A_n(\theta,\phi)\,\mathrm{d}\Omega},
  \qquad G = \eta\,D,
  \label{eq:directivity}
\end{equation}
where:
\begin{itemize}[leftmargin=1.5em,itemsep=1pt,topsep=2pt]
  \item $A_n(\theta,\phi)$ is the normalised power pattern (equal to 1 at
        boresight), as a function of the usual spherical angles
        $(\theta,\phi)$;
  \item $\eta$ is the radiation efficiency (fraction of input power actually
        radiated);
  \item $\mathrm{d}\Omega$ is an element of solid angle, in steradians \citep{balanis2016,stutzman2013}.
\end{itemize}
The \emph{half-power beamwidth} (HPBW) is the angular width over which
$A_n \geq 0.5$, and it sets the field of view of a single element. Once the
sky position is expressed in the direction-cosine coordinates $(l,m)$ used
later in this chapter, the same normalised pattern is written $A_n(l,m)$; in
an interferometer this is exactly the \emph{primary beam} that multiplies the
sky brightness in Section~\ref{sec:visibility-to-image}. Using one
isolated-element pattern for every antenna in the array assumes that the
elements do not interact electromagnetically with one another; this
assumption is checked directly, by a mutual-coupling study of the real array,
in Chapter~\ref{ch:antenna}.

For imaging purposes the key quantity derived from the gain is the
\emph{effective aperture} $A_e(\theta,\phi)$,
\begin{equation}
  A_e = \frac{\lambda^2}{4\pi}\,G ,
  \label{eq:aeff}
\end{equation}
where $\lambda$ is the observing wavelength and $G$ is the gain of
Eq.~\eqref{eq:directivity} \citep{balanis2016,tms2017}; $A_e$ has units of
area and is, loosely, the area over which the antenna effectively collects
power from an incident plane wave.

The receiving element used throughout this thesis is a Yagi-Uda array: a
single driven element (here a folded dipole), a reflector behind it, and
several directors in front. The parasitic elements re-radiate with phases set
by their length and spacing, producing an end-fire pattern with gain typically
in the range \SIrange{8}{12}{dBi}. The folded dipole raises the driven-element
impedance to roughly four times that of a simple dipole, which eases matching
to the feed. Besides the gain and HPBW already defined above, the relevant
figure of merit for how well the element is matched to its feed line is the
standing-wave ratio
\begin{equation}
  \mathrm{SWR} = \frac{1+|\Gamma|}{1-|\Gamma|},
  \qquad
  \Gamma = \frac{Z_\mathrm{in}-Z_0}{Z_\mathrm{in}+Z_0},
  \label{eq:swr}
\end{equation}
where $Z_\mathrm{in}$ is the antenna input impedance and $Z_0$ is the
feed-line impedance (\SI{75}{\ohm} for the coaxial cable used)
\citep{pozar2012,balanis2016}. These quantities are computed for the real
element in Chapter~\ref{ch:antenna}.

\subsection{Friis Reception Relation}
\label{subsec:friis}
With the effective aperture available, a source of flux density $S_\nu$
observed over a bandwidth $\Delta\nu$ delivers, to a matched load, a received
power given by the Friis reception relation~\citep{friis1946}:
\begin{equation}
  P_\mathrm{rx} = \tfrac{1}{2}\,A_e\,S_\nu\,\Delta\nu ,
  \label{eq:friis}
\end{equation}
where:
\begin{itemize}[leftmargin=1.5em,itemsep=1pt,topsep=2pt]
  \item $S_\nu$ is the flux density of the source, in
        $\si{\watt\per\square\metre\per\hertz}$ (Section~\ref{subsec:flux-sfu});
  \item $\Delta\nu$ is the receiver bandwidth, in hertz;
  \item $A_e$ is the effective aperture of Eq.~\eqref{eq:aeff};
  \item the factor $\tfrac12$ accounts for a single received polarisation of
        unpolarised radiation \citep{tms2017,kraus1986};
  \item $P_\mathrm{rx}$ is the resulting received power, in watts.
\end{itemize}
Equation~\eqref{eq:friis} is the bridge between the astronomical flux scale
and the receiver, and it is a computational step evaluated directly by the
simulation: it takes $S_\nu(t)$ from Section~\ref{subsec:sv-of-t} and turns it
into a received power $P_\mathrm{rx}(t)$, as implemented in
Chapter~\ref{ch:pipeline}.

\subsection{Voltage Equivalent at the Receiver Input}
\label{subsec:voltage-equiv}
The simulation, like a real digitiser, works with voltages, not power. For a
receiver with a resistive load $R_\mathrm{load}$, the received power of
Eq.~\eqref{eq:friis} corresponds to a root-mean-square signal voltage
\begin{equation}
  V_\mathrm{rms} = \sqrt{P_\mathrm{rx}\,R_\mathrm{load}} ,
  \label{eq:vrms}
\end{equation}
where $P_\mathrm{rx}$ is the received power of Eq.~\eqref{eq:friis} and
$R_\mathrm{load}$ is the load resistance seen by the antenna; the relation
follows directly from $P=V_\mathrm{rms}^2/R$ for a matched resistive load.
Equation~\eqref{eq:vrms} is the last step of the chain from flux to voltage,
and, like Eq.~\eqref{eq:friis}, it is evaluated directly by the simulation in
Chapter~\ref{ch:pipeline} (Section~\ref{sec:voltage-synth}). The antenna's
directional response enters at this same voltage level, by multiplying
$V_\mathrm{rms}$ with the square root of the normalised primary beam $A_n$
evaluated at the Sun's instantaneous position, so that a source away from
boresight produces a smaller voltage than the same source at boresight. This
gives the deterministic, signal part of the antenna voltage; the next section
adds the part that is not deterministic.

% =====================================================================
\section{Receiver Noise and Averaging}
\label{sec:noise}
Even though the Sun produces a measurable flux density, a real receiver never
observes that signal alone. Every channel also carries thermal noise
contributed by the antenna itself, the sky, the ground, the cable, the
low-noise amplifier, and the rest of the receiver chain. A faithful
simulation must therefore generate not only the deterministic solar voltage of
Eq.~\eqref{eq:vrms}, but also a statistically correct noise voltage to add to
it. This section builds that noise model and then asks how averaging can
recover a weak signal from underneath it.

\subsection{System Temperature and Johnson-Nyquist Noise}
\label{subsec:johnson-noise}
Any resistive element at physical temperature $T$ generates a fluctuating
voltage, even with no signal present~\citep{nyquist1928}. The available noise
power in a bandwidth $\Delta\nu$ is, in the Rayleigh-Jeans limit,
\begin{equation}
  P_n = k_B\,T\,\Delta\nu ,
  \label{eq:johnson}
\end{equation}
where:
\begin{itemize}[leftmargin=1.5em,itemsep=1pt,topsep=2pt]
  \item $k_B$ is Boltzmann's constant, as in Eq.~\eqref{eq:rj};
  \item $T$ is the physical temperature of the noise-generating element, in
        kelvin;
  \item $\Delta\nu$ is the bandwidth, in hertz;
  \item $P_n$ is the resulting noise power, in watts.
\end{itemize}
Equation~\eqref{eq:johnson} underlies the convention of expressing every
noise contribution as an equivalent \emph{noise temperature}. The total
\emph{system temperature} $T_\mathrm{sys}$ sums the sky, ground, antenna, and
receiver contributions, and the system-equivalent flux density follows from
Eq.~\eqref{eq:friis} as $\mathrm{SEFD}=2k_B T_\mathrm{sys}/A_e$
\citep{tms2017,kraus1986}. Exactly as for the signal in
Section~\ref{subsec:voltage-equiv}, the simulation needs a voltage, not a
power: the noise power of Eq.~\eqref{eq:johnson} is converted to a noise
voltage with standard deviation
\begin{equation}
  \sigma_n = \sqrt{\tfrac{1}{2}\,P_n\,R_\mathrm{load}} ,
  \label{eq:noise-voltage}
\end{equation}
where the factor $\tfrac12$ splits the total noise power equally between the
in-phase and quadrature components of the complex baseband signal. In the
simulation, $T_\mathrm{sys}$ is a single configurable parameter, and
Eqs.~\eqref{eq:johnson}-\eqref{eq:noise-voltage} are applied independently to
every antenna to draw a fresh noise voltage at every time step, exactly as
independent thermal noise would arise in every physical receiver channel.
Adding this noise voltage to the signal voltage of Eq.~\eqref{eq:vrms} gives
the complete per-antenna complex voltage used in Chapter~\ref{ch:pipeline}
(Section~\ref{sec:voltage-synth}).

\subsection{Radiometer Equation, and Why Averaging Helps}
\label{subsec:radiometer}
Equation~\eqref{eq:noise-voltage} shows that noise is always present, but it
does not yet say whether a given signal can be recovered from underneath it.
The \emph{radiometer equation} answers that question: it gives the smallest
detectable change in antenna temperature after integrating for a time $\tau$
over a bandwidth $\Delta\nu$,
\begin{equation}
  \Delta T_\mathrm{rms} = \frac{T_\mathrm{sys}}{\sqrt{\Delta\nu\,\tau}} ,
  \label{eq:radiometer}
\end{equation}
where $T_\mathrm{sys}$ is the system temperature of
Section~\ref{subsec:johnson-noise}, $\Delta\nu$ is the bandwidth, and $\tau$
is the integration (averaging) time \citep{tms2017,kraus1986}. This equation
is not itself a step that generates simulated data; its role is to explain
\emph{why} the time-averaging performed by the correlator in
Chapter~\ref{ch:pipeline} improves the signal-to-noise ratio: thermal noise
is a random process with zero mean, so summing more independent samples,
either by integrating longer or by using a wider bandwidth, averages the
noise toward zero while the steady solar signal adds coherently. This
$1/\sqrt{\Delta\nu\,\tau}$ scaling is the statistical foundation on which both
single-dish detection and interferometric correlation rest, and it is
reproduced explicitly by the noise model of Section~\ref{sec:voltage-synth} in
Chapter~\ref{ch:pipeline}.

\subsection{Signal-to-Noise Interpretation for the Simulation}
\label{subsec:snr-interpretation}
Combining Eqs.~\eqref{eq:friis}, \eqref{eq:johnson}, and \eqref{eq:radiometer}
gives a direct way to ask whether a given signal is detectable: the quiet-Sun
or burst flux density is converted to an equivalent antenna temperature and
compared with $\Delta T_\mathrm{rms}$ for the integration time and bandwidth
in use. If the equivalent signal temperature is well above $\Delta
T_\mathrm{rms}$, the source should be visible above the noise floor in the
simulated voltages and, after correlation, in the simulated visibilities; if
it is comparable to or below $\Delta T_\mathrm{rms}$, no amount of correct
geometry will recover a clean fringe, which is exactly the situation
encountered in the field deployment described in Chapter~\ref{ch:intro}. This
comparison is carried out quantitatively for the simulated instrument in
Section~\ref{sec:snr-analysis} of Chapter~\ref{ch:results}.

% =====================================================================
\section{From Antenna Voltage to Interferometric Visibility}
\label{sec:voltage-to-visibility}
Every antenna in the array now has its own complex voltage: a deterministic
part from Eq.~\eqref{eq:vrms} and a random noise part from
Eq.~\eqref{eq:noise-voltage}. A single antenna, however, can only measure
\emph{power}-the mean value of the squared voltage-and power alone carries no
information about \emph{where} on the sky the signal came from. Recovering
that positional information is the job of interferometry, and it starts from
a simple geometric fact: a celestial source is composed of a vast number of
independent radiating regions-electrons gyrating incoherently in the corona,
for instance-so the radiation received from any two distinct directions on
the sky is statistically independent. This property, that the source is
\emph{spatially incoherent}, is the single most important physical assumption
in interferometry: the coherence measured \emph{between two antennas} does
not arise because the source itself is coherent, but because the same
incoherent field is sampled at two separated points, and the geometric delay
between those points encodes information about the source's angular
structure. This is precisely the principle exploited in the early twentieth
century by the Michelson stellar interferometer, in which two mirror
apertures, fed into a common eyepiece, produce fringes whose contrast falls as
the angular size of the source increases relative to $\lambda/b$; it is the
optical-astronomy ancestor of the radio-interferometric technique developed
below.

\subsection{Geometric Delay and Phase}
\label{subsec:geometric-delay}
Consider two identical antennas, separated by a baseline vector $\bm{b}$, that
both receive the wavefront from a point source in direction $\hat{\bm{s}}$.
Because the two antennas are not at the same point in space, the wavefront
reaches one of them slightly later than the other. This time difference is
the \emph{geometric delay} \citep{taylor1999synthesis,tms2017}
\begin{equation}
  \tau_g = \frac{\bm{b}\cdot\hat{\bm{s}}}{c},
  \label{eq:taug}
\end{equation}
where $\bm{b}$ is the baseline vector between the two antennas, $\hat{\bm{s}}$
is the unit vector pointing toward the source, and $c$ is the speed of light.
At a centre frequency $\nu_0$, this delay corresponds to a phase difference
$2\pi\nu_0\tau_g$ between the two antenna voltages: this is the entire
physical origin of everything that follows in this section. Because the
Earth rotates, $\hat{\bm{s}}$ changes slowly with time relative to $\bm{b}$,
so $\tau_g$, and therefore the phase difference, also changes with time. In a
real instrument the bulk of this delay is compensated electronically-a
process called \emph{fringe stopping}-by inserting a compensating delay
$\tau_g$ in one signal path before correlation, so that the residual delay
tracks only the difference between the true geometric delay and the model
delay used by the correlator.

\subsection{Two-Element Correlation}
\label{subsec:two-element-correlation}
A \emph{correlator} multiplies the voltage outputs of two antennas and
time-averages the product. For a single point source, the correlator output
after fringe stopping and averaging over many carrier cycles is the real part
of a complex response:
\begin{equation}
  r(\tau_g) \propto \cos\!\big(2\pi\nu_0\,\tau_g\big)
            = \cos\!\Big(\frac{2\pi}{\lambda}\,\bm{b}\cdot\hat{\bm{s}}\Big).
  \label{eq:fringe}
\end{equation}
This is the celebrated \emph{interference fringe} \citep{taylor1999synthesis}:
as the Earth rotates, $\hat{\bm{s}}$ sweeps slowly across the sky relative to
$\bm{b}$, and the correlator output traces out a slowly varying cosine whose
instantaneous frequency is the \emph{fringe rate}. In practice a \emph{complex
correlator} forms both an in-phase output ($\cos$) and a quadrature output
($\sin$, obtained by inserting a quarter-wave phase shift in one path before
the second multiplication), so that the full complex response, amplitude and
phase, of the source on that baseline is recovered, rather than only its real
part.

For an east-west baseline of length $b$ observing a source at hour angle $H$
and declination $\delta$, differentiating Eq.~\eqref{eq:taug} with respect to
time gives the rate of change of fringe phase, and hence the fringe period
\begin{equation}
  T_\mathrm{fringe} \approx \frac{\lambda}{b\,\omega_\oplus\,\cos\delta},
  \label{eq:fringeperiod}
\end{equation}
where $\lambda$ is the observing wavelength, $b$ is the baseline length, $\delta$
is the source declination, and $\omega_\oplus$ is the Earth's sidereal
rotation rate \citep{taylor1999synthesis,tms2017}. For $b=\SI{8}{\metre}$,
$\lambda=\SI{0.49}{\metre}$ and the solar declination at the spring equinox
($\delta = 0^\circ$), this evaluates to approximately $14$~minutes.
Equation~\eqref{eq:fringeperiod} is \emph{not} used to generate the simulated
data; it plays a different role, as an independent analytic prediction used
only afterwards, in Chapter~\ref{ch:results}, to check whether the simulated
fringe behaves the way theory says it must. It is the most basic and most
directly verifiable analytic signature of an interferometer, and for that
reason it is the first validation benchmark used in this thesis.

The derivation above implicitly assumed a single frequency $\nu_0$. In
practice every receiver has finite bandwidth $\Delta\nu$, and correlating over
that bandwidth attenuates the fringe amplitude if the residual delay error
after fringe stopping, $\Delta\tau = \tau_g - \tau_\mathrm{model}$, is not
negligible. Integrating a uniform passband of width $\Delta\nu$ centred on
$\nu_0$ introduces a multiplicative attenuation factor
\begin{equation}
  \mathrm{sinc}\!\big(\Delta\nu\,\Delta\tau\big)
  = \frac{\sin(\pi\,\Delta\nu\,\Delta\tau)}{\pi\,\Delta\nu\,\Delta\tau},
  \label{eq:bandwidth-smearing}
\end{equation}
known as \emph{bandwidth smearing} \citep{tms2017}. Because the residual delay
error in this simulation is determined entirely by the known geometry rather
than by an imperfect delay model, this effect is small for the narrow field of
the quiet Sun and is quantified in Chapter~\ref{ch:pipeline}.

\subsection{Visibility Amplitude and Phase}
\label{subsec:visibility-amp-phase}
Section~\ref{subsec:two-element-correlation} described the correlator in
terms of an idealised continuous field. In terms of the actual per-antenna
complex voltages built in Sections~\ref{sec:flux-to-voltage}
and~\ref{sec:noise}, $v_i(t)$ and $v_j(t)$, what the correlator for antenna
pair $(i,j)$ computes is
\begin{equation}
  V_{ij} = \big\langle v_i(t)\,v_j^{*}(t) \big\rangle ,
  \label{eq:correlator-def}
\end{equation}
where:
\begin{itemize}[leftmargin=1.5em,itemsep=1pt,topsep=2pt]
  \item $v_i(t)$ is the complex voltage at antenna $i$;
  \item $v_j^{*}(t)$ is the complex conjugate of the voltage at antenna $j$;
  \item $\langle\cdot\rangle$ denotes a time average, long enough to average
        down the noise (Section~\ref{subsec:radiometer}) but short enough
        that the source itself is effectively unchanged.
\end{itemize}
The result $V_{ij}$ is a single complex number that combines both an
amplitude and a phase: this number is the \emph{complex visibility},
$V=|V|\,e^{i\phi}$. Its amplitude $|V|$ measures the degree of coherence of
the field between the two antennas-equal to the total received flux for a
baseline short enough to be insensitive to the source's angular size, and
falling as the baseline lengthens and begins to resolve the source-while its
phase $\phi$ encodes the position of the source relative to the
phase-tracking centre. This is the key conceptual difference between a single
antenna and an interferometer: \emph{a single antenna measures power; a pair
of antennas measures spatial coherence}, and it is this spatial-coherence
measurement, repeated over many baselines and over time, that ultimately
allows an image to be reconstructed.

For an extended, spatially incoherent source, Eq.~\eqref{eq:correlator-def}
generalises the point-source fringe of Eq.~\eqref{eq:fringe} to an integral
over the whole sky brightness distribution $I(\hat{\bm{s}})$:
\begin{equation}
  V(\bm{b}) = \int_{\Omega} I(\hat{\bm{s}})\,
  \exp\!\Big[-i\,2\pi\nu_0\,\frac{\bm{b}\cdot\hat{\bm{s}}}{c}\Big]\,\mathrm{d}\Omega ,
  \label{eq:coherence-relation}
\end{equation}
where the integral is taken over the solid angle $\Omega$ occupied by the
source \citep{taylor1999synthesis,tms2017}. Equation~\eqref{eq:coherence-relation}
already contains the essential physics of synthesis imaging: the visibility
measured on a baseline $\bm{b}$ is a weighted, phase-shifted integral of the
sky brightness, so an array that samples many different baselines is, in
effect, sampling many different weighted integrals of the same unknown sky.
Turning this statement into a precise, invertible imaging equation requires a
convenient coordinate system for the baseline and the sky, which
Section~\ref{sec:uvw-coverage} introduces next, and the formal statement of
the imaging relation is given in Section~\ref{sec:visibility-to-image}.

% =====================================================================
\section{From Physical Baselines to \texorpdfstring{$uvw$}{uvw} Coverage}
\label{sec:uvw-coverage}
Equation~\eqref{eq:coherence-relation} is written in terms of a baseline
vector $\bm{b}$ and a source direction $\hat{\bm{s}}$, but it says nothing yet
about how $\bm{b}$ is actually obtained from the physical layout of the array,
or how it changes as the source moves across the sky. That is the purpose of
this section: it connects the antenna positions that are physically surveyed
in the field to the coordinate frame in which the visibility integral is most
naturally expressed.

\subsection{ENU Ground Coordinates}
\label{subsec:enu}
Every antenna in the array has a fixed physical position on the ground,
measured in a local East-North-Up (ENU) frame tangent to the Earth at the
observing site. The baseline vector $\bm{b}$ used throughout
Section~\ref{sec:voltage-to-visibility} is simply the difference between the
ENU positions of two antennas. This is the natural frame in which an array is
laid out and surveyed in the field, and it is the starting point of the
coordinate chain used in Chapter~\ref{ch:pipeline} to describe an arbitrary
two-dimensional array layout, not only a single east-west line.

\subsection{Source-Based \texorpdfstring{$uvw$}{uvw} Coordinates}
\label{subsec:uvw-transform}
Visibility, however, is most naturally expressed in a frame tied to the
source direction, not to the ground: this is why Eq.~\eqref{eq:coherence-relation}
is written with $\bm{b}\cdot\hat{\bm{s}}$ rather than with the raw ENU
coordinates. The array geometry must therefore be transformed from
ground-based ENU coordinates to source-based spatial-frequency coordinates,
conventionally called $(u,v,w)$. Obtaining that frame takes two steps, as
described by \citet{taylor1999synthesis}. First, the local ENU antenna
positions are converted to an Earth-fixed equatorial $(X,Y,Z)$ frame, in which
$X$ points to the intersection of the equator and the local meridian, $Y$
points east, and $Z$ points to the celestial pole. Second, this equatorial
frame is rotated, for the instantaneous hour angle $H$ and declination
$\delta$ of the phase-tracking centre, into the baseline-dependent $(u,v,w)$
frame in which $w$ points along the source direction:
\begin{equation}
  \begin{bmatrix} u\\ v\\ w \end{bmatrix}
  =
  \begin{bmatrix}
    \sin H & \cos H & 0\\
    -\sin\delta\cos H & \sin\delta\sin H & \cos\delta\\
    \cos\delta\cos H & -\cos\delta\sin H & \sin\delta
  \end{bmatrix}
  \begin{bmatrix} X_\lambda\\ Y_\lambda\\ Z_\lambda \end{bmatrix},
  \label{eq:uvw}
\end{equation}
where:
\begin{itemize}[leftmargin=1.5em,itemsep=1pt,topsep=2pt]
  \item $H$ is the hour angle of the phase-tracking centre;
  \item $\delta$ is its declination;
  \item $(X_\lambda,Y_\lambda,Z_\lambda)$ are the equatorial baseline
        components, expressed in units of wavelength;
  \item $(u,v,w)$ are the resulting source-based baseline coordinates, also
        in units of wavelength, with $u$ and $v$ spanning the plane
        perpendicular to the source direction and $w$ along it
        \citep{taylor1999synthesis,tms2017}.
\end{itemize}
This single, general transformation is applied identically to every baseline
and every instant in the simulation pipeline, so that per-antenna geometric
phase, instantaneous baseline length, and fringe-stopping delay are all
computed from one consistent geometric model, valid for arbitrary
two-dimensional array layouts rather than only for the east-west case of
Eq.~\eqref{eq:taug}.

\subsection{Earth Rotation and \texorpdfstring{$uv$}{uv}-Coverage}
\label{subsec:earth-rotation-uv}
The set of $(u,v)$ points sampled by an interferometric array-its
\emph{$uv$-coverage}-dictates the array's imaging capability, in direct
analogy to the way the aperture of a filled-aperture telescope determines its
point-spread function. A single snapshot with $N$ antennas yields $N(N-1)/2$
instantaneous baselines, hence $N(N-1)/2$ independent samples of $V(u,v)$
(each baseline contributing two complex-conjugate samples at $\pm(u,v)$, since
$I(l,m)$ is real and therefore $V(-u,-v)=V^{*}(u,v)$). For the four-element
array of this thesis, $N=4$ gives only six instantaneous baselines-far too
sparse, on its own, to constrain a two-dimensional image.

As the Earth rotates, however, the projection of any fixed physical baseline
onto the plane perpendicular to the source direction changes continuously,
because the rotation matrix in Eq.~\eqref{eq:uvw} depends on the hour angle
$H$. For a baseline with equatorial components $(X_\lambda,Y_\lambda,Z_\lambda)$
observing a source at fixed declination $\delta$, eliminating $H$ between the
$u$ and $v$ equations of Eq.~\eqref{eq:uvw} shows that each baseline traces an
\emph{ellipse} in the $(u,v)$ plane as $H$ varies over the day, with
semi-major axis equal to the projected east-west baseline length and
semi-minor axis foreshortened by $\sin\delta$. This technique, known as
\emph{Earth-rotation synthesis} \citep{taylor1999synthesis}, allows a small
number of physical antennas to synthesise, over the course of a transit, a
$uv$-coverage equivalent to a much larger number of simultaneous baselines,
without any moving parts. It is the reason a four-element, six-baseline array
can nonetheless produce a meaningful synthesised image once data accumulated
over an observing transit are combined-a trade-off explored quantitatively,
for the compact east-west geometry of the physical array against more general
two-dimensional configurations, in Chapter~\ref{ch:results}.

Earth-rotation synthesis assumes that each visibility sample is attributed to
a single, exact $(u,v)$ point, when in fact it is an average over the
integration time of one correlator dump. If the integration time is long
enough for $H$ to change appreciably, a source away from the phase-tracking
centre is smeared tangentially in the $uv$ plane, an effect known as
\emph{time-average smearing}, in direct analogy to the bandwidth smearing of
Eq.~\eqref{eq:bandwidth-smearing}. For the short integration times and narrow
field of the quiet Sun used in this thesis this effect is small, and it is
quantified together with bandwidth smearing in Chapter~\ref{ch:pipeline}.

% =====================================================================
\section{From Visibility to Dirty Image}
\label{sec:visibility-to-image}

\subsection{Van Cittert-Zernike Theorem}
\label{subsec:vcz}
Visibility is useful precisely because, under the two assumptions already
introduced-an incoherent source (Section~\ref{sec:voltage-to-visibility}) and
a narrow field of view-it is related to the sky brightness by an ordinary
two-dimensional Fourier transform. In this small-field limit, appropriate for
the compact solar disk observed in this thesis, the relation is
\begin{equation}
  V(u,v) = \iint A_n(l,m)\,I(l,m)\,e^{-2\pi i(ul+vm)}\,\mathrm{d}l\,\mathrm{d}m ,
  \label{eq:vcz}
\end{equation}
where:
\begin{itemize}[leftmargin=1.5em,itemsep=1pt,topsep=2pt]
  \item $I(l,m)$ is the sky brightness distribution;
  \item $A_n(l,m)$ is the normalised primary beam of a single antenna
        (Section~\ref{subsec:antenna-gain-beam});
  \item $(l,m)$ are direction cosines on the sky, measured relative to the
        phase-tracking centre;
  \item $(u,v)$ are the baseline components in wavelengths, from
        Section~\ref{sec:uvw-coverage}.
\end{itemize}
This is the \emph{van Cittert-Zernike theorem} as used throughout synthesis
imaging \citep{taylor1999synthesis,tms2017}: the visibility and the (apodised)
sky brightness are Fourier conjugates, so that, given sufficient and suitable
sampling of $V(u,v)$, the sky brightness can in principle be recovered by
inverse Fourier transformation. Each baseline therefore samples one point in
this Fourier $(u,v)$ plane, and it is the combination of many simultaneous
baselines with the additional points that Earth rotation supplies over time
(Section~\ref{subsec:earth-rotation-uv}) that builds up usable $uv$-coverage.

Equation~\eqref{eq:vcz} is a simplification of a more general relation that
keeps a third baseline coordinate, $w$, and a corresponding phase term:
\begin{equation}
  V(u,v,w) = \iint A_n(l,m)\,I(l,m)\,
  e^{-2\pi i\,[\,ul + vm + w(n-1)\,]}\,
  \frac{\mathrm{d}l\,\mathrm{d}m}{n},
  \qquad n=\sqrt{1-l^2-m^2} .
  \label{eq:visibility}
\end{equation}
If the primary beam is narrow enough that $l,m\ll1$ over the region where
$A_nI\neq0$, then $n-1\approx-\tfrac12(l^2+m^2)$ is small, and provided the
$w$-term phase error $2\pi w(n-1)$ remains $\ll1$ radian across the field of
view, the $w$-dependence may be dropped entirely and Eq.~\eqref{eq:visibility}
reduces exactly to Eq.~\eqref{eq:vcz}. For the solar disk, whose
\SI{32}{\arcminute} angular diameter subtends $l,m\sim5\times10^{-3}$, and for
the short, near-zenith baselines used in this thesis, the $w$-term is indeed
negligible, and the small-field relation of Eq.~\eqref{eq:vcz} applies without
correction. For a wider field of view or a longer-baseline array than used in
this thesis, the $w$-term of Eq.~\eqref{eq:visibility} would have to be
retained.

A baseline of length $|\bm{b}|$ measured in wavelengths corresponds, through
Eq.~\eqref{eq:vcz}, to an angular scale of order $1/|\bm{b}|_\lambda$ radians
on the sky. For a uniformly bright circular disk of angular radius
$\theta_s$ (a simplistic but effective model for the quiet Sun), the
two-dimensional Fourier transform in Eq.~\eqref{eq:vcz} has the closed
analytic form
\begin{equation}
  V(q) = \frac{2 J_1(2\pi q \theta_s)}{2\pi q \theta_s} \equiv \jinc(2\pi q\theta_s),
  \qquad q=\sqrt{u^2+v^2},
  \label{eq:jinc}
\end{equation}
where $J_1$ is the first-order Bessel function of the first kind and $q$ is
the baseline length in wavelengths \citep{tms2017,taylor1999synthesis}. The
$\jinc$ function falls monotonically from unity at $q=0$ (the disk's total
flux, measured on a zero-length baseline) to its first null at
$q\theta_s\approx0.61$, beyond which the visibility amplitude is small and
oscillatory. Like the fringe period of Eq.~\eqref{eq:fringeperiod},
Eq.~\eqref{eq:jinc} is not used to generate simulated data; it is an
independent analytic target, computed separately from the simulation, against
which the simulated visibilities of the resolved solar disk are validated in
Chapter~\ref{ch:results}: a correctly implemented simulation must reproduce
both the $\jinc$ envelope and the location of its first null as a function of
baseline length.

Two further practical error sources act on the phase of the visibility rather
than on its amplitude, and are worth noting here because they affect the
imaging step directly. \emph{Phase-calibration and pointing errors} shift and
blur the reconstructed image; in this simulation the geometry is known
exactly, so they are introduced only as controllable perturbations
(Chapter~\ref{ch:pipeline}). \emph{Ionospheric phase} is, for the short
baselines used in this thesis ($\lesssim\SI{130}{\metre}$ at
\SI{610}{\mega\hertz}), essentially common-mode between antennas and cancels
in the differential (baseline) measurement, with a residual differential
phase of about one degree.

\subsection{Direct Fourier Inversion}
\label{subsec:fourier-inversion}
Because only a finite, discrete set of $(u,v)$ points is ever sampled, it is
convenient to describe the measurement formally as a multiplication of the
true, continuous visibility function by a \emph{sampling function} $S(u,v)$-a
sum of Dirac delta functions (optionally weighted to account for the relative
sensitivity or density of different baselines) located at each measured
$(u,v)$ point and its conjugate. The measured visibility is then
$S(u,v)\,V(u,v)$, and inverse-Fourier-transforming this sampled, rather than
continuous, visibility does not yield the true sky brightness $I$, but rather
the \emph{dirty image} $I_\mathrm{D}$:
\begin{equation}
  I_\mathrm{D}(l,m) = \iint S(u,v)\,V(u,v)\,e^{+2\pi i(ul+vm)}\,\mathrm{d}u\,\mathrm{d}v .
  \label{eq:dirty-ft}
\end{equation}
where $S(u,v)$ is the sampling function defined above and $V(u,v)$ is the
measured visibility of Eq.~\eqref{eq:vcz} at each sampled point.
Equation~\eqref{eq:dirty-ft} is the direct inversion that the simulation
actually carries out in Chapter~\ref{ch:pipeline}: every measured visibility
is placed at its $(u,v)$ location and the sampled function is
inverse-Fourier-transformed, with no further processing such as
deconvolution.

\subsection{Dirty Beam, PSF, and the Convolution Relation}
\label{subsec:dirty-beam}
By the convolution theorem \citep{taylor1999synthesis}, multiplication of $V$
by $S$ in the Fourier ($uv$) domain corresponds to convolution in the image
domain, so Eq.~\eqref{eq:dirty-ft} may equivalently, and more usefully, be
written as the \emph{Convolution Equation} of synthesis imaging
\citep{taylor1999synthesis}:
\begin{equation}
  I_\mathrm{D}(l,m) = I(l,m) * B(l,m),
  \label{eq:dirty}
\end{equation}
where $*$ denotes two-dimensional convolution, $I(l,m)$ is the true sky
brightness, and the \emph{dirty beam} (equivalently, the
\emph{point-spread function}, PSF, or \emph{synthesised beam}) $B(l,m)$ is
simply the Fourier transform of the sampling function itself,
\begin{equation}
  B(l,m) = \iint S(u,v)\,e^{+2\pi i(ul+vm)}\,\mathrm{d}u\,\mathrm{d}v .
  \label{eq:dirtybeam}
\end{equation}
Equations~\eqref{eq:dirty}-\eqref{eq:dirtybeam} \citep{taylor1999synthesis,tms2017}
express the fundamental imaging equation of an interferometer: the dirty
image is exactly what one would expect from imaging the true sky through an
instrument whose response to a point source is the beam $B(l,m)$-a direct
radio analogue of the diffraction-limited PSF of an optical telescope, except
that here the PSF is determined not by a filled aperture but by the discrete
pattern of sampled $(u,v)$ points. Put practically: the dirty image is
\emph{not} a perfect picture of the true sky. It is the true sky "stamped"
by the array's PSF, and because $S(u,v)$ is a sparse, typically non-uniform
set of samples rather than a filled, azimuthally symmetric disk (as for a
filled-aperture telescope), the dirty beam $B(l,m)$ generally has a narrow
central peak-whose width sets the nominal angular resolution of the
synthesised image, of order $1/b_\mathrm{max}$ radians for a maximum baseline
$b_\mathrm{max}$ measured in wavelengths-surrounded by sidelobes that reflect
the gaps in $uv$-coverage. For the compact, six-baseline east-west array of
this thesis, the $uv$-coverage accumulated over an Earth-rotation transit
consists of six partial elliptical arcs rather than a densely filled disk, so
$B(l,m)$ retains significant sidelobe structure; this is expected, and the
resulting dirty image is still a rigorously correct image of the sky under
Eq.~\eqref{eq:dirty} given the actual sampling achieved-it is simply not a
clean, sidelobe-free representation of $I(l,m)$. Because sidelobes are an
unavoidable consequence of sparse $uv$-coverage rather than a simulation
defect, Chapter~\ref{ch:results} evaluates the simulated dirty image not only
by inspection but through the PSF shape itself, its dynamic range, and its
main-lobe FWHM.

Because the simulation in this thesis observes a single, compact, time-stable
calibrator (the quiet-Sun disk) rather than attempting to image a complex
field crowded with confusing sources, the dirty image and dirty beam together
are sufficient to validate the instrument: demonstrating numerically that the
simulated $I_\mathrm{D}$ equals the assumed input sky map convolved with the
independently computed $B(l,m)$, as required by Eq.~\eqref{eq:dirty},
constitutes the ultimate end-to-end mathematical verification of the entire
signal-processing pipeline, from antenna voltages to final image, and is
carried out explicitly in Chapter~\ref{ch:results}.

% =====================================================================
\section{Chapter Summary: Theoretical Chain Used by the Simulation}
\label{sec:chapter-summary}
Read as a whole, this chapter tells one story. The Sun produces radio
emission (Section~\ref{sec:solar-activity}); the international community
measures this emission as a flux density, in solar flux units
(Section~\ref{sec:radioastro-quantities}); an antenna converts that flux
density into received radio-frequency power, and a receiver converts that
power into a voltage together with its own unavoidable noise
(Sections~\ref{sec:flux-to-voltage}-\ref{sec:noise}); several antennas at
known positions see the same signal with a small, geometry-dependent phase
difference, and a correlator turns a pair of antenna voltages into a single
complex visibility (Section~\ref{sec:voltage-to-visibility}); the array
geometry places every visibility at its correct point in a Fourier $(u,v)$
plane (Section~\ref{sec:uvw-coverage}); and an inverse Fourier transform of
those points produces the dirty image analysed in Chapter~\ref{ch:results}
(Section~\ref{sec:visibility-to-image}). Every equation in this chapter is a
necessary step in that one chain, not an independent piece of theory.

Not every equation plays the same role in the chain, and it is worth
separating the two roles explicitly, as flagged throughout this chapter:
\begin{itemize}[leftmargin=1.5em]
  \item \textbf{Computational steps, evaluated directly by the simulation in
        Chapter~\ref{ch:pipeline}:} the flux-to-SI conversion
        (Eq.~\eqref{eq:sfu}) and interpolation to $S_\nu(t)$
        (Section~\ref{subsec:sv-of-t}); the Friis relation and voltage
        equivalent (Eqs.~\eqref{eq:friis}-\eqref{eq:vrms}); the Johnson-Nyquist
        noise power and noise voltage (Eqs.~\eqref{eq:johnson}-\eqref{eq:noise-voltage});
        the $\mathrm{ENU}\to XYZ\to uvw$ geometry (Eq.~\eqref{eq:uvw}); and
        the direct Fourier inversion to the dirty image
        (Eq.~\eqref{eq:dirty-ft}).
  \item \textbf{Analytic benchmarks, used only afterwards to validate the
        simulation in Chapter~\ref{ch:results}:} the fringe period
        (Eq.~\eqref{eq:fringeperiod}); the van Cittert-Zernike relation and the
        closed-form $\jinc$ visibility of a uniform disk
        (Eqs.~\eqref{eq:vcz}-\eqref{eq:jinc}); and the convolution relation
        between the dirty image, the true sky, and the synthesised beam
        (Eqs.~\eqref{eq:dirty}-\eqref{eq:dirtybeam}).
\end{itemize}
This formalism, developed throughout in the notation of
\citet{taylor1999synthesis}, provides both the modelling choices used to build
the simulation and the analytic benchmarks used to check it, and it is this
distinction-not merely the equations themselves-that the chapters that follow
rely on.
